\chapter{Vecteurs, translations et combinaisons}
\section*{Exercice 8.1 -- Déplacements}
$\overrightarrow{AB}=(3,4)$. On ajoute ce déplacement à $C$, donc $D(-1,6)$.

\section*{Exercice 8.2 -- Combinaison}
$3\vec u-2\vec v=(6,-9)-(-2,8)=(8,-17)$. Sa norme vaut $\sqrt{64+289}=\sqrt{353}$.

\section*{Exercice 8.3 -- Chasles}
Les deux premiers vecteurs donnent $\overrightarrow{AC}$. Puis $-\overrightarrow{DC}=\overrightarrow{CD}$, donc la somme vaut $\overrightarrow{AC}+\overrightarrow{CD}=\overrightarrow{AD}$.

\section*{Exercice 8.4 -- Parallélogramme}
L'égalité $\overrightarrow{AD}=\overrightarrow{BC}=(1,4)$ donne $D(2,5)$. Les diagonales ont alors le milieu commun $(3,3{,}5)$, ce qui fournit une vérification.

\section*{Exercice 8.5 -- Milieu vectoriel}
Par Chasles, la somme vaut $2\overrightarrow{MI}+\overrightarrow{IA}+\overrightarrow{IB}$. Les deux derniers vecteurs sont opposés puisque $I$ est le milieu. Il reste $2\overrightarrow{MI}$.

\section*{Exercice 8.6 -- Changement d'échelle}
La direction est conservée, le sens inversé et la norme vaut $(3/2)\times5=7{,}5$. Si $\vec u=(a,b)$, la nouvelle norme est $\sqrt{(9/4)(a^2+b^2)}=(3/2)\sqrt{a^2+b^2}$.
